% =====================================================================
%  Derivations, conventions and diagnostics
%  Companion to "Dissolution of the collinear structure in the
%  solar-sail circular restricted three-body problem".
%  Not for submission: this is the long-form working document, with
%  every algebraic step written out and every residual explained.
% =====================================================================
\documentclass[11pt,a4paper]{article}

\usepackage[margin=1in]{geometry}
\usepackage{amsmath,amssymb,bm}
\usepackage{booktabs}
\usepackage{siunitx}
\usepackage[numbers]{natbib}
\usepackage{hyperref}
\usepackage[font=small,labelfont=bf]{caption}

\hypersetup{
  colorlinks = true,
  linkcolor  = {[rgb]{0.12,0.22,0.45}},
  citecolor  = {[rgb]{0.10,0.36,0.26}},
  urlcolor   = {[rgb]{0.40,0.16,0.36}},
  pdftitle   = {Derivations, conventions and diagnostics},
  pdfauthor  = {Bishwaswarup}
}

\sisetup{group-minimum-digits=4, detect-all}

\newcommand{\rhat}{\hat{\bm{r}}}
\newcommand{\phat}{\hat{\bm{p}}}
\newcommand{\qhat}{\hat{\bm{q}}}
\newcommand{\nhat}{\hat{\bm{n}}}

\title{\bfseries Derivations, conventions and diagnostics\\
\large A companion to \emph{Dissolution of the collinear structure in the
solar-sail circular restricted three-body problem}}
\author{Bishwaswarup}
\date{\today}

\begin{document}
\maketitle

\begin{abstract}
\noindent
This is the long-form working document behind the preprint. Nothing here is new
physics; everything is the algebra the paper compresses or states without proof,
written out step by step, together with an account of what every reported
residual actually measures. It exists so that each number in the manuscript can
be checked by hand rather than taken on trust. Section~\ref{sec:resid} in
particular explains the floating-point quantities --- the $10^{-16}$'s --- which
are easy to misread as physics.

Two honest caveats are recorded here rather than buried: the
\num{8.1e-10} quoted in the manuscript's frequency scan is an artefact of grid
alignment and overstates that scan's resolving power
(Section~\ref{sec:scanresid}), and the $\Delta_z$ branch discriminant weakens by
an order of magnitude across the sail band (Section~\ref{sec:deltaz}).
\end{abstract}

\tableofcontents

% =====================================================================
\section{Conventions, units and notation}\label{sec:conv}
% =====================================================================

\subsection{What is set to unity}

The circular restricted three-body problem is non-dimensionalised so that three
independent quantities are unity:
\begin{enumerate}
\item the separation of the two primaries, $L=1$;
\item the total mass, $M_1+M_2=1$, so that the primaries have masses $1-\mu$
      and $\mu$ with $\mu=M_2/(M_1+M_2)$;
\item the mean motion of the primaries about their common barycentre, $n=1$,
      which with $L=1$ and $G(M_1+M_2)=1$ is Kepler's third law
      $n^2L^3=G(M_1+M_2)$.
\end{enumerate}
Setting $n=1$ fixes the unit of time. One non-dimensional (nd) time unit is
$1/n$, and one full revolution of the primaries is $2\pi$ nd. Hence
\begin{equation}
  1\ \text{nd time unit} = \frac{T_{\mathrm{orbit}}}{2\pi}
  = \frac{\SI{365.25636}{\day}}{2\pi} = \SI{58.1324}{\day}
  \label{eq:timeunit}
\end{equation}
for Sun--Earth. This conversion is used nowhere in the manuscript except in the
instability-timescale section, and it is the only place a physical time enters.

Lengths convert with the primary separation. For Sun--Earth that is one
astronomical unit,
\begin{equation}
  1\ \text{nd length} = \SI{1.495978707e8}{\kilo\metre} .
  \label{eq:lengthunit}
\end{equation}

\subsection{Geometry and notation}

The synodic (co-rotating) frame places the larger primary at $x=-\mu$ and the
smaller at $x=1-\mu$, both on the $x$-axis, with the $z$-axis normal to the
orbital plane. Distances from the primaries are
\begin{equation}
  r_1^2 = (x+\mu)^2+y^2+z^2, \qquad
  r_2^2 = (x-1+\mu)^2+y^2+z^2 .
\end{equation}

Table~\ref{tab:notation} collects the symbols. The single most important one is
$\beta$, the \emph{lightness number}: the ratio of the sail's
radiation-pressure acceleration to the local solar gravitational acceleration.
Because both fall off as $1/r_1^2$, this ratio is a constant of the sail, not a
function of position --- which is exactly why a face-on sail can be absorbed
into a rescaling of the primary's mass.

\begin{table}[t]
  \centering
  \caption{Notation.}
  \label{tab:notation}
  \begin{tabular}{ll}
    \toprule
    symbol & meaning \\
    \midrule
    $\mu$                    & mass ratio $M_2/(M_1+M_2)$ \\
    $\beta$                  & lightness number, sail acceleration $/$ solar gravity \\
    $r_1,\;r_2$              & distances to the larger and smaller primary \\
    $\gamma$                 & local length scale $\lvert(1-\mu)-x_{\mathrm{eq}}\rvert$ \\
    $r_H$                    & Hill radius $(\mu/3)^{1/3}$ \\
    $U$                      & modified pseudo-potential, eq.~\eqref{eq:U} \\
    $A\equiv c_2$            & second Legendre coefficient of the collinear expansion \\
    $S,\;s$                  & the solar and tidal contributions to $A$, eq.~\eqref{eq:Asplit} \\
    $\kappa$                 & the level at which $s$ is held, $s=\kappa$ \\
    $\lambda_u$              & unstable (saddle) exponent, in-plane \\
    $\omega$                 & in-plane centre frequency \\
    $\nu$                    & out-of-plane frequency \\
    $\tau_u$                 & instability e-folding time, $1/\lambda_u$ \\
    $\alpha,\;\delta$        & sail cone and clock angles \\
    \bottomrule
  \end{tabular}
\end{table}

% =====================================================================
\section{Equations of motion}\label{sec:eom}
% =====================================================================

\subsection{The rotating frame}

In an inertial frame the massless particle obeys Newtonian gravitation from two
point masses on circular orbits. Transforming to the frame rotating with the
primaries at unit angular velocity $\bm{\Omega}=\hat{\bm{z}}$ introduces the
centrifugal and Coriolis accelerations,
\begin{equation}
  \bm{a}_{\mathrm{rot}} = \bm{a}_{\mathrm{in}}
  - 2\,\bm{\Omega}\times\dot{\bm{r}}
  - \bm{\Omega}\times(\bm{\Omega}\times\bm{r}) .
\end{equation}
With $\bm{\Omega}=\hat{\bm{z}}$ one has
$\bm{\Omega}\times\dot{\bm{r}}=(-\dot y,\dot x,0)$ and
$-\bm{\Omega}\times(\bm{\Omega}\times\bm{r})=(x,y,0)$. The centrifugal term is
the gradient of $\tfrac12(x^2+y^2)$, so it can be folded into a potential; the
Coriolis term cannot, because it depends on velocity. This asymmetry is the
reason the in-plane and out-of-plane linearisations behave so differently
(Section~\ref{sec:lin}) and, ultimately, the reason a purely transverse control
input can still reach the in-plane radial direction (Section~\ref{sec:control}).

Collecting terms, and writing $P$ and $Q$ for the two gravitational strengths,
\begin{equation}
  \ddot x - 2\dot y = \frac{\partial U}{\partial x}, \qquad
  \ddot y + 2\dot x = \frac{\partial U}{\partial y}, \qquad
  \ddot z         = \frac{\partial U}{\partial z},
  \label{eq:eom}
\end{equation}
\begin{equation}
  U = \tfrac12(x^2+y^2) + \frac{P}{r_1} + \frac{Q}{r_2},
  \qquad P=(1-\beta)(1-\mu), \quad Q=\mu .
  \label{eq:U}
\end{equation}
Setting $\beta=0$ gives $P=1-\mu$ and recovers the classical CR3BP.

\subsection{The sail force and why \texorpdfstring{$P=(1-\beta)(1-\mu)$}{P=(1-beta)(1-mu)}}

For an ideal specular reflector the force is along the sail normal with
magnitude proportional to $\cos^2\alpha$:
\begin{equation}
  \bm{a}_{\mathrm{sail}} = \beta\,\frac{1-\mu}{r_1^2}\,\cos^2\alpha\;\nhat ,
  \qquad
  \nhat = \cos\alpha\,\rhat + \sin\alpha\cos\delta\,\phat
        + \sin\alpha\sin\delta\,\qhat .
  \label{eq:sail}
\end{equation}
One factor of $\cos\alpha$ comes from the projected area the sail presents to
the Sun; the second comes from the fact that only the component of incoming
momentum normal to the sail is reversed. A perfect \emph{absorber} differs: the
force lies along the Sun line with magnitude $\propto\cos\alpha$, one factor
only. The two models coincide at $\alpha=0$, which is the only case the bulk of
the manuscript uses.

At $\alpha=0$, $\nhat=\rhat$ and \eqref{eq:sail} reduces to
\begin{equation}
  \bm{a}_{\mathrm{sail}} = \beta\,\frac{1-\mu}{r_1^2}\,\rhat ,
\end{equation}
which is a central inverse-square force directed away from the larger primary.
It is therefore derivable from the potential $-\beta(1-\mu)/r_1$, and adding it
to the gravitational term $+(1-\mu)/r_1$ gives $+(1-\beta)(1-\mu)/r_1$. This is
the whole content of the statement that \emph{a face-on sail rescales solar
gravity}: the problem is not perturbed, it is a CR3BP with a different mass for
the larger primary. Every exact result in the manuscript follows from this one
observation.

For $\alpha\neq0$ the force acquires a transverse component, is no longer
central, and no potential exists --- hence no Jacobi integral, and hence the
restriction of Sections 3--6 of the manuscript to $\alpha=0$.

\subsection{The Jacobi integral}

Because \eqref{eq:eom} is autonomous with a velocity-independent potential and a
Coriolis term that does no work,
\begin{equation}
  \frac{\mathrm{d}}{\mathrm{d}t}\Bigl(2U-v^2\Bigr)
  = 2\nabla U\cdot\dot{\bm r} - 2\,\dot{\bm r}\cdot\ddot{\bm r} = 0
\end{equation}
after substituting \eqref{eq:eom}, since
$\dot{\bm r}\cdot(\bm\Omega\times\dot{\bm r})=0$. Therefore
\begin{equation}
  C_{\mathrm{sail}} = 2U - v^2
  \label{eq:jacobi}
\end{equation}
is conserved, with $U$ the \emph{modified} potential \eqref{eq:U}. The
gravitational Jacobi constant built from $(1-\mu)/r_1$ instead differs by
\begin{equation}
  C_{\mathrm{grav}} = C_{\mathrm{sail}} + \frac{2\beta(1-\mu)}{r_1},
  \label{eq:cgrav}
\end{equation}
and since $r_1$ varies along an orbit, $C_{\mathrm{grav}}$ is \emph{not}
conserved for $\beta\neq0$. Section~\ref{sec:jacobiresid} quantifies both.

% =====================================================================
\section{The exact identity \texorpdfstring{$r_1+r_2=1$}{r1+r2=1}}\label{sec:unitsum}
% =====================================================================

On the segment of the $x$-axis strictly between the primaries,
$-\mu<x<1-\mu$, one has $x+\mu>0$ and $x-1+\mu<0$, so
\begin{equation}
  r_1 = x+\mu, \qquad r_2 = (1-\mu)-x,
  \label{eq:r1r2}
\end{equation}
and therefore
\begin{equation}
  r_1+r_2 = (x+\mu) + (1-\mu) - x = 1
  \label{eq:sum}
\end{equation}
identically, for every $x$ in the interval and every $\mu$. This is not an
approximation and involves no expansion in $\mu$. It is worth emphasising
because the standard treatment of collinear points works with
$\gamma=r_2$ and expands in $\gamma$, generating series in $\mu^{1/3}$; here no
series is needed. Equation~\eqref{eq:sum} is what makes the closed form of
Section~\ref{sec:closed} exact rather than asymptotic: imposing a condition on
$r_2$ immediately fixes $r_1=1-r_2$ with no error term.

% =====================================================================
\section{Linearisation at a collinear equilibrium}\label{sec:lin}
% =====================================================================

\subsection{The equilibrium condition}

An equilibrium on the axis has $y=z=\dot x=\dot y=\dot z=0$ and
$\partial U/\partial x=0$. Differentiating \eqref{eq:U},
\begin{equation}
  \frac{\partial U}{\partial x}
  = x - \frac{P(x+\mu)}{r_1^3} - \frac{Q(x-1+\mu)}{r_2^3} .
\end{equation}
Using \eqref{eq:r1r2}, $x+\mu=r_1$ and $x-1+\mu=-r_2$, so
\begin{equation}
  \frac{\partial U}{\partial x} = x - \frac{P}{r_1^2} + \frac{Q}{r_2^2},
\end{equation}
and the equilibrium satisfies
\begin{equation}
  \boxed{\;x = \frac{P}{r_1^2} - \frac{Q}{r_2^2}
           = \frac{(1-\beta)(1-\mu)}{r_1^2} - \frac{\mu}{r_2^2}\;}
  \label{eq:balance}
\end{equation}
This is the on-axis force balance used throughout. Read physically: the
centrifugal term $x$ balances the reduced solar attraction minus the smaller
primary's pull.

\subsection{Second derivatives on the axis}

Write $\Phi=1/r_1$. Then
\begin{equation}
  \frac{\partial\Phi}{\partial x} = -\frac{x+\mu}{r_1^3},
  \qquad
  \frac{\partial^2\Phi}{\partial x^2}
  = -\frac{1}{r_1^3} + \frac{3(x+\mu)^2}{r_1^5} .
\end{equation}
Evaluated on the axis, $(x+\mu)^2=r_1^2$, so
\begin{equation}
  \left.\frac{\partial^2\Phi}{\partial x^2}\right|_{\mathrm{axis}}
  = -\frac{1}{r_1^3} + \frac{3r_1^2}{r_1^5} = \frac{2}{r_1^3} .
\end{equation}
For the transverse directions, $\partial^2\Phi/\partial y^2
= -1/r_1^3 + 3y^2/r_1^5$, which on the axis ($y=0$) is $-1/r_1^3$; identically
for $z$. The same computation for $1/r_2$ gives $+2/r_2^3$ and $-1/r_2^3$,
because $(x-1+\mu)^2=r_2^2$ on the axis regardless of sign.

Adding the centrifugal part $\tfrac12(x^2+y^2)$, whose second derivatives are
$1,1,0$:
\begin{equation}
  U_{xx} = 1 + 2A, \qquad U_{yy} = 1 - A, \qquad U_{zz} = -A,
  \label{eq:hessian}
\end{equation}
\begin{equation}
  A = \underbrace{\frac{(1-\beta)(1-\mu)}{r_1^3}}_{\textstyle S}
    + \underbrace{\frac{\mu}{r_2^3}}_{\textstyle s} .
  \label{eq:Asplit}
\end{equation}
All the cross terms $U_{xy},U_{xz},U_{yz}$ vanish on the axis by symmetry, so
the Hessian is diagonal and the $z$ motion decouples exactly at linear order.

\subsection{The out-of-plane mode}

From \eqref{eq:eom} and \eqref{eq:hessian}, the vertical perturbation
$\zeta=z$ obeys
\begin{equation}
  \ddot\zeta = U_{zz}\,\zeta = -A\,\zeta
  \qquad\Longrightarrow\qquad
  \ddot\zeta + A\zeta = 0 ,
\end{equation}
a pure harmonic oscillator with
\begin{equation}
  \boxed{\;\nu = \sqrt{A}\;}
  \label{eq:nu}
\end{equation}
provided $A>0$, which holds for every collinear point. There is no Coriolis
coupling out of plane because $\bm\Omega\times\dot{\bm r}$ has no $z$
component. \textbf{Physical meaning:} $\nu$ is the frequency at which a
spacecraft displaced above the ecliptic oscillates through it. Its period is
$2\pi/\nu$ in nd time.

\subsection{The in-plane modes and the characteristic quartic}

Write $\xi=x-x_{\mathrm{eq}}$, $\eta=y$. The linearised in-plane system is
\begin{equation}
  \ddot\xi - 2\dot\eta = (1+2A)\,\xi, \qquad
  \ddot\eta + 2\dot\xi = (1-A)\,\eta .
  \label{eq:inplane}
\end{equation}
Seek $\xi=X\mathrm{e}^{\lambda t}$, $\eta=Y\mathrm{e}^{\lambda t}$:
\begin{align}
  (\lambda^2 - 1 - 2A)\,X - 2\lambda\,Y &= 0, \\
  2\lambda\,X + (\lambda^2 - 1 + A)\,Y &= 0 .
\end{align}
A non-trivial solution requires the determinant to vanish:
\begin{equation}
  (\lambda^2-1-2A)(\lambda^2-1+A) + 4\lambda^2 = 0 .
\end{equation}
Put $u=\lambda^2$ and expand fully:
\begin{align}
  (u-1-2A)(u-1+A)
  &= u^2 - u + Au - u + 1 - A - 2Au + 2A - 2A^2 \notag\\
  &= u^2 + u(-2 - A) + (1 + A - 2A^2) .
\end{align}
Adding $4u$,
\begin{equation}
  \boxed{\;u^2 + (2-A)\,u + \bigl(1+A-2A^2\bigr) = 0\;}
  \label{eq:quartic}
\end{equation}
which is the quartic $\lambda^4+(2-A)\lambda^2+(1+A-2A^2)=0$ in $\lambda$.

\subsection{Solving it}

The discriminant is
\begin{align}
  (2-A)^2 - 4(1+A-2A^2)
  &= 4 - 4A + A^2 - 4 - 4A + 8A^2 \notag\\
  &= 9A^2 - 8A = A(9A-8) \;\equiv\; D^2 .
\end{align}
Hence
\begin{equation}
  u_{\pm} = \tfrac12\Bigl[(A-2) \pm D\Bigr], \qquad
  D = \sqrt{A(9A-8)} .
  \label{eq:roots}
\end{equation}
The product of the roots is $u_+u_- = 1+A-2A^2 = -(2A+1)(A-1)$. For $A>1$ this
is strictly negative, so the two roots have opposite signs: one positive, one
negative. Therefore
\begin{equation}
  \boxed{\;\lambda_u^2 = \tfrac12\bigl[(A-2)+D\bigr], \qquad
          \omega^2 = \tfrac12\bigl[D-(A-2)\bigr]\;}
  \label{eq:rates}
\end{equation}
with $\lambda_u$ real (the saddle) and $\omega$ real (the centre, since
$u_-=-\omega^2<0$ gives $\lambda=\pm\mathrm{i}\omega$).

The condition $A>1$ is exactly the condition for a real hyperbolic pair. At
$A=1$ the product vanishes and one root goes to zero; for $A<1$ both roots are
negative and the point is linearly stable in plane. Every classical collinear
point has $A>1$, and a sail drives $A\to1^+$ from above but never below
(Section~\ref{sec:closed}).

\subsection{What the three rates mean}\label{sec:meaning}

This is worth stating plainly, because the manuscript uses all three without
defining them operationally.

\begin{description}
\item[$\lambda_u$ --- the saddle exponent.] A perturbation along the unstable
eigenvector grows as $\mathrm{e}^{\lambda_u t}$. Its reciprocal
$\tau_u=1/\lambda_u$ is the \emph{e-folding time}: the time for a departure to
grow by a factor $\mathrm{e}\approx2.718$. This is the number that determines
station-keeping cadence. For the classical Sun--Earth $L_1$,
$\lambda_u=\num{2.53256}$ nd, so by \eqref{eq:timeunit}
$\tau_u=\SI{58.1324}{\day}/2.53256=\SI{22.95}{\day}$ --- the familiar figure
quoted for $L_1$ halo missions. That the model returns it without being tuned
to is a check on \eqref{eq:rates}.

\item[$\omega$ --- the in-plane centre frequency.] The planar Lyapunov orbits
about the equilibrium have period $2\pi/\omega$ at infinitesimal amplitude.
Physically it is the epicyclic frequency of the libration in the
$(x,y)$ plane, modified by the two primaries' tidal fields.

\item[$\nu$ --- the out-of-plane frequency.] The vertical Lyapunov orbits have
period $2\pi/\nu$ at infinitesimal amplitude.
\end{description}

\noindent\textbf{Why the ratio $\nu/\omega$ matters.} A halo orbit is a
three-dimensional periodic orbit in which the in-plane and out-of-plane motions
share a common period. At linear order this is impossible unless
$\nu/\omega=1$ exactly. Halo families therefore bifurcate from the planar
Lyapunov family at finite amplitude, where nonlinear frequency corrections close
the gap. The closer $\nu/\omega$ is to unity at linear order, the smaller the
amplitude at which the halo family appears. The manuscript's bound
\eqref{eq:band} says this gap is never larger than \SI{5.72}{\percent} and
never zero.

% =====================================================================
\section{The identity \texorpdfstring{$A\equiv c_2$}{A = c2}}\label{sec:c2}
% =====================================================================

Richardson's expansion of the collinear dynamics writes the potential as a
series in Legendre polynomials with coefficients
\begin{equation}
  c_n = \frac{1}{\gamma^3}\left[\mu
        + (-1)^n\,\frac{(1-\mu)\,\gamma^{\,n+1}}{(1-\gamma)^{\,n+1}}\right],
\end{equation}
$\gamma$ being the distance from the smaller primary to the libration point.
For $n=2$, with $\gamma=r_2$ and $1-\gamma=r_1$ by \eqref{eq:sum},
\begin{equation}
  c_2 = \frac{\mu}{r_2^3} + \frac{1-\mu}{r_1^3},
\end{equation}
which is precisely \eqref{eq:Asplit} at $\beta=0$. Hence $A$ is not a new
quantity: it is the second Legendre coefficient, and the sail's sole effect is
to replace $(1-\mu)$ by $(1-\beta)(1-\mu)$ inside it. This is the compact
statement that \emph{a sail continuously detunes $c_2$}, a parameter otherwise
fixed once the system and the libration point are chosen. The manuscript
verifies the identity numerically to $<10^{-13}$ at three systems.

% =====================================================================
\section{The tidal-parity threshold, in full}\label{sec:closed}
% =====================================================================

\subsection{The condition}

Define the \emph{saddle strength} $s=\mu/r_2^3$, the smaller primary's direct
contribution to $A$ in \eqref{eq:Asplit}. Tidal parity is the level
\begin{equation}
  s = 1
  \qquad\Longleftrightarrow\qquad
  \frac{\mu}{r_2^3}=1
  \qquad\Longleftrightarrow\qquad
  r_2 = \mu^{1/3} .
  \label{eq:parity}
\end{equation}
Since the Hill radius is $r_H=(\mu/3)^{1/3}$,
\begin{equation}
  \frac{r_2}{r_H} = \frac{\mu^{1/3}}{(\mu/3)^{1/3}} = 3^{1/3} = 1.4422\ldots
\end{equation}
independently of $\mu$ --- the parity radius is always $3^{1/3}$ Hill radii.

\subsection{The substitution, term by term}

Impose \eqref{eq:parity} on the balance \eqref{eq:balance}. Three substitutions
are needed, and each is exact:
\begin{align}
  r_2 &= \mu^{1/3}, \\
  r_1 &= 1 - r_2 = 1 - \mu^{1/3}
        &&\text{by \eqref{eq:sum}, no error term}, \\
  \frac{\mu}{r_2^2} &= \frac{\mu}{\mu^{2/3}} = \mu^{1/3}
        &&\text{the tidal term equals the standoff}, \\
  x &= (1-\mu) - r_2 = (1-\mu) - \mu^{1/3}
        &&\text{by \eqref{eq:r1r2}} .
\end{align}
The third line is the crux and deserves emphasis: at $s=1$, and only there, the
smaller primary's gravitational acceleration at the equilibrium,
$\mu/r_2^2$, is numerically equal to the equilibrium's distance from that
primary, $r_2$. Substituting into \eqref{eq:balance},
\begin{equation}
  \underbrace{(1-\mu)-\mu^{1/3}}_{x}
  = \frac{(1-\beta)(1-\mu)}{\bigl(1-\mu^{1/3}\bigr)^2}
    - \underbrace{\mu^{1/3}}_{\mu/r_2^2} .
\end{equation}
The $-\mu^{1/3}$ appears on both sides and cancels identically, leaving
\begin{equation}
  1-\mu = \frac{(1-\beta)(1-\mu)}{\bigl(1-\mu^{1/3}\bigr)^2} .
\end{equation}
Dividing by $(1-\mu)$, which is non-zero for $\mu<1$,
\begin{equation}
  1-\beta = \bigl(1-\mu^{1/3}\bigr)^2
\end{equation}
and therefore
\begin{equation}
  \boxed{\;\beta_{\mathrm{crit}}
   = 1 - \bigl(1-\mu^{1/3}\bigr)^2
   = 2\mu^{1/3} - \mu^{2/3}\;}
  \label{eq:closedform}
\end{equation}
Expanding the square gives $1-(1-2\mu^{1/3}+\mu^{2/3})=2\mu^{1/3}-\mu^{2/3}$.

\subsection{Why this is exact}

Three features make it so. First, \eqref{eq:sum} is an identity, not a leading
term. Second, the condition $s=1$ fixes $r_2$ in closed form, so nothing needs
to be solved numerically. Third, the cancellation is algebraic --- it does not
rely on $\mu$ being small. Consequently \eqref{eq:closedform} holds for
$\mu=\tfrac12$ as well as for $\mu=10^{-7}$. To leading order
$\beta_{\mathrm{crit}}\simeq2\mu^{1/3}$, but the exact form costs nothing.

\subsection{The linear structure at parity}

At parity $r_2=\mu^{1/3}$ and $r_1=1-\mu^{1/3}$, so from \eqref{eq:Asplit}
\begin{equation}
  S = \frac{(1-\beta)(1-\mu)}{r_1^3}
    = \frac{(1-\mu^{1/3})^2(1-\mu)}{(1-\mu^{1/3})^3}
    = \frac{1-\mu}{1-\mu^{1/3}} ,
\end{equation}
using $1-\beta=r_1^2$. Hence
\begin{equation}
  A^{\mathrm{parity}} = 1 + \frac{1-\mu}{1-\mu^{1/3}} ,
  \label{eq:Aparity}
\end{equation}
which for Sun--Earth is $2.014635$, not $2$. The naive identification
$A-1=s$ is therefore short by
\begin{equation}
  \frac{A-2}{A} = \frac{0.014635}{2.014635} = \SI{0.7265}{\percent},
\end{equation}
quoted as \SI{0.73}{\percent} in the manuscript. The discrepancy is
$S-1=O(\mu^{1/3})$ and reflects the smaller primary's \emph{indirect} effect:
it displaces the equilibrium, which changes $r_1$, which changes $S$.

% =====================================================================
\section{The tidal-strength hierarchy}\label{sec:kappa}
% =====================================================================

\subsection{Generalising the level}

Replace \eqref{eq:parity} by $s=\kappa$ for arbitrary $\kappa>0$:
\begin{equation}
  \frac{\mu}{r_2^3}=\kappa
  \quad\Longrightarrow\quad
  r_2(\kappa) = \Bigl(\frac{\mu}{\kappa}\Bigr)^{1/3}
              = \Bigl(\frac{3}{\kappa}\Bigr)^{1/3} r_H .
  \label{eq:r2kappa}
\end{equation}
Again $r_1=1-r_2$ exactly. The tidal term is now
\begin{equation}
  \frac{\mu}{r_2^2} = \mu\Bigl(\frac{\kappa}{\mu}\Bigr)^{2/3}
                    = \mu^{1/3}\kappa^{2/3},
\end{equation}
while the standoff is $r_2=\mu^{1/3}\kappa^{-1/3}$. Substituting into
\eqref{eq:balance},
\begin{equation}
  (1-\mu)-r_2 = \frac{(1-\beta)(1-\mu)}{r_1^2} - \mu^{1/3}\kappa^{2/3},
\end{equation}
so
\begin{equation}
  \frac{(1-\beta)(1-\mu)}{r_1^2}
  = (1-\mu) \;\underbrace{-\,r_2 + \mu^{1/3}\kappa^{2/3}}_{\textstyle\Delta(\kappa)} .
\end{equation}
The residue is
\begin{equation}
  \Delta(\kappa) = -\mu^{1/3}\kappa^{-1/3} + \mu^{1/3}\kappa^{2/3}
                 = \mu^{1/3}\kappa^{-1/3}\bigl(\kappa-1\bigr),
  \label{eq:residue}
\end{equation}
and rearranging gives the general closed form
\begin{equation}
  \boxed{\;\beta(\kappa) = 1 - \frac{(1-r_2)^2}{1-\mu}
    \Bigl[(1-\mu) + \mu^{1/3}\kappa^{-1/3}(\kappa-1)\Bigr]\;}
  \label{eq:betakappa}
\end{equation}

\subsection{Why \texorpdfstring{$\kappa=1$}{kappa=1} is distinguished}

From \eqref{eq:residue}, $\Delta(\kappa)=0$ if and only if $\kappa=1$, since
$\mu^{1/3}\kappa^{-1/3}>0$ for all admissible $\mu,\kappa$. At that level and
no other, \eqref{eq:betakappa} collapses to
\begin{equation}
  1-\beta = (1-r_2)^2 = \bigl(1-\mu^{1/3}\bigr)^2,
\end{equation}
recovering \eqref{eq:closedform}.

The statement to make is therefore precise, and slightly weaker than one might
hope: a closed form exists for \emph{every} $\kappa$ --- the algebra never
requires a numerical root. What is special about $\kappa=1$ is that the smaller
primary's direct term drops out of the balance identically, leaving a threshold
that is a pure square in $\mu^{1/3}$ with no residual $\kappa$-dependence. This
answers the objection that $s=1$ is merely a choice of units: it is the unique
level at which the tidal acceleration and the standoff coincide,
$\mu/r_2^2=r_2$, and that coincidence is a property of the force balance.

\subsection{Both of the manuscript's markers are members}

Because $r_2/r_H=(3/\kappa)^{1/3}$ by \eqref{eq:r2kappa}:
\begin{itemize}
\item $\kappa=3$ gives $r_2/r_H=1$ exactly --- the Hill-sphere exit criterion,
      at $\beta=\num{2.98e-4}$;
\item $\kappa=1$ gives $r_2/r_H=3^{1/3}=1.4422$ --- tidal parity, at
      $\beta=\num{0.028646}$;
\item the classical $L_1$ itself sits at $\kappa=s(\beta{=}0)=3.0303$, just
      inside the Hill sphere at $r_2=0.9967\,r_H$.
\end{itemize}
The two markers the manuscript treats separately --- one adopted, one dismissed
as ``not a useful discriminant'' --- are the same object at two levels.

% =====================================================================
\section{The frequency ratio}\label{sec:freq}
% =====================================================================

\subsection{Forming the ratio}

From \eqref{eq:nu} and \eqref{eq:rates}, with $c_2\equiv A$,
\begin{equation}
  \frac{\nu^2}{\omega^2}
  = \frac{c_2}{\tfrac12\bigl[D-(c_2-2)\bigr]}
  = \frac{2c_2}{D-(c_2-2)},
  \qquad D=\sqrt{c_2(9c_2-8)} .
  \label{eq:R}
\end{equation}
Note this depends on $c_2$ alone --- $\mu$ and $\beta$ enter only through
$c_2$. Any statement about this ratio is therefore a statement about the
collinear normal form itself, shared by every three-body system.

\subsection{The stationary point}

Write $R=2c/M$ with $M=D-c+2$ and $D=\sqrt{9c^2-8c}$. Then
\begin{equation}
  \frac{\mathrm{d}D}{\mathrm{d}c}=\frac{18c-8}{2D}=\frac{9c-4}{D},
  \qquad
  M' = \frac{9c-4}{D}-1 .
\end{equation}
By the quotient rule,
\begin{equation}
  \frac{\mathrm{d}R}{\mathrm{d}c} = \frac{2M-2cM'}{M^2},
\end{equation}
which vanishes when $M=cM'$:
\begin{equation}
  D-c+2 = c\left(\frac{9c-4}{D}-1\right) = \frac{c(9c-4)}{D}-c .
\end{equation}
The $-c$ cancels on both sides, leaving $D+2=c(9c-4)/D$, i.e.
\begin{equation}
  D^2+2D = 9c^2-4c .
\end{equation}
Since $D^2=9c^2-8c$,
\begin{equation}
  9c^2-8c+2D = 9c^2-4c
  \quad\Longrightarrow\quad
  2D = 4c
  \quad\Longrightarrow\quad
  D = 2c .
\end{equation}
Squaring, $4c^2=9c^2-8c$, so $5c^2=8c$ and (discarding $c=0$)
\begin{equation}
  \boxed{\;c_2^\star = \tfrac85\;}
  \label{eq:cstar}
\end{equation}

\subsection{The rates there, and the \texorpdfstring{$7:8:9$}{7:8:9} ratio}

At $c_2=8/5$, $D=2c_2=16/5$. From \eqref{eq:rates} and \eqref{eq:nu}:
\begin{align}
  \lambda_u^2 &= \tfrac12\bigl[(\tfrac85-2)+\tfrac{16}5\bigr]
               = \tfrac12\bigl[-\tfrac25+\tfrac{16}5\bigr]
               = \tfrac12\cdot\tfrac{14}5 = \tfrac75, \\
  \nu^2       &= c_2 = \tfrac85, \\
  \omega^2    &= \tfrac12\bigl[\tfrac{16}5+\tfrac25\bigr]
               = \tfrac12\cdot\tfrac{18}5 = \tfrac95 .
\end{align}
Hence
\begin{equation}
  \lambda_u^2:\nu^2:\omega^2 = \tfrac75:\tfrac85:\tfrac95 = 7:8:9,
\end{equation}
and
\begin{equation}
  \left.\frac{\nu}{\omega}\right|_{c_2^\star}
  = \sqrt{\tfrac89} = \frac{2\sqrt2}{3} = 0.942809041582\ldots
  \label{eq:bound}
\end{equation}

\subsection{The endpoints}

\textbf{As $c_2\to1^+$:} $D=\sqrt{1\cdot1}=1$, so
$\omega^2=\tfrac12[1-(-1)]=1$ and $\nu^2=1$; the ratio tends to $1$. This is
the hovering limit, where the planar roots become $\lambda^2\in\{0,-1\}$ and
the in-plane period is driven to exactly $2\pi$ --- one year for Sun--Earth.
That is the Keplerian epicyclic degeneracy of an inverse-square field, not a
resonance with the smaller primary.

\textbf{As $c_2\to\infty$:} $D=3c_2\sqrt{1-8/(9c_2)}\simeq3c_2-\tfrac43$, so
\begin{equation}
  \omega^2 \simeq \tfrac12\bigl[3c_2-\tfrac43-c_2+2\bigr] = c_2+\tfrac13,
  \qquad
  \frac{\nu^2}{\omega^2}\simeq\frac{c_2}{c_2+\tfrac13}\to1 .
\end{equation}
So the ratio approaches unity at both ends and dips to $2\sqrt2/3$ in between.
The band is therefore
\begin{equation}
  \frac{\nu}{\omega}\in\Bigl[\tfrac{2\sqrt2}{3},\,1\Bigr),
  \qquad \text{width } 1-\tfrac{2\sqrt2}{3}=\SI{5.719}{\percent} .
  \label{eq:band}
\end{equation}
The interval is closed on the left (the bound is attained at $c_2=8/5$) and
open on the right (unity is approached but never reached for finite $c_2>1$).

% =====================================================================
\section{Which commensurabilities the band admits}\label{sec:reso}
% =====================================================================

The manuscript originally tested fourteen low-order rationals and found none
inside \eqref{eq:band}. That is true but weaker than necessary, and the
stronger statement is elementary.

Let $p/q$ with $0<p<q$ be a rational in lowest terms. For a fixed denominator
$q$, the largest such rational strictly below $1$ is $(q-1)/q$. Therefore
$p/q\in[2\sqrt2/3,1)$ requires
\begin{equation}
  \frac{q-1}{q}\;\geq\;\frac{2\sqrt2}{3}
  \quad\Longleftrightarrow\quad
  1-\frac1q\geq\frac{2\sqrt2}{3}
  \quad\Longleftrightarrow\quad
  q \geq \frac{1}{1-2\sqrt2/3} = 17.4853\ldots
\end{equation}
so $q\geq18$. Hence:
\begin{itemize}
\item \textbf{No} rational with denominator $q\leq17$ lies in the band at all.
\item The lowest-denominator one that does is $17{:}18=0.9444\ldots$, of
      order $p+q=35$.
\item The fourteen tested rationals have $q\leq9$ and are excluded without
      being tested individually.
\end{itemize}

It is important to state the converse honestly: the band is \emph{not}
resonance-free. It is an interval of positive length, so it contains infinitely
many rationals --- $18{:}19$, $19{:}20$, and so on. The correct claim is that
every commensurability it admits is of order $35$ or higher, which is far beyond
what any perturbation theory would retain, and that no commensurability
underwrites tidal parity.

% =====================================================================
\section{The instability timescale}\label{sec:tau}
% =====================================================================

Define
\begin{equation}
  \tau_u = \frac{1}{\lambda_u}, \qquad
  N_{\mathrm{efold}} = \lambda_u\,T_{\mathrm{mission}} ,
  \label{eq:tau}
\end{equation}
with $\lambda_u$ in nd units. Converting with \eqref{eq:timeunit},
\begin{equation}
  \tau_u\,[\mathrm{days}] = \frac{\SI{58.1324}{\day}}{\lambda_u\,[\mathrm{nd}]} .
\end{equation}
At $\beta=0$, $\lambda_u=2.53256$ gives $\tau_u=\SI{22.95}{\day}$. At tidal
parity, $\lambda_u=1.505418$ gives $\tau_u=\SI{38.62}{\day}$.

Inverting $N_{\mathrm{efold}}=1$ for $\beta$ requires a numerical root because
$\lambda_u(\beta)$ is transcendental: one solves
$\lambda_u(\beta)=\SI{58.1324}{\day}/T$ by bisection on the monotone decreasing
function $\lambda_u$. The results are in the manuscript's Table 4. Note this
threshold is \emph{not} $\mu$-independent and \emph{not} exact --- it depends on
the mission horizon chosen, which is an engineering input, not a property of the
dynamics. It is reported because it is the operationally meaningful reading,
not because it is a sharper result than \eqref{eq:closedform}.

% =====================================================================
\section{Sail control authority}\label{sec:control}
% =====================================================================

\subsection{The orthonormal triad}

The triad in \eqref{eq:sail} must be orthonormal. Take
\begin{equation}
  \rhat = \frac{\bm r-\bm r_1}{r_1}, \qquad
  \phat = \frac{\hat{\bm k}\times\rhat}{\lVert\hat{\bm k}\times\rhat\rVert},
  \qquad
  \qhat = \rhat\times\phat .
\end{equation}
A common shortcut substitutes $\hat{\bm k}$ itself for $\qhat$. That fails off
the ecliptic because $\rhat\cdot\hat{\bm k}=z/r_1\neq0$, and then
\begin{equation}
  \lVert\nhat\rVert^2 = 1 + \sin(2\alpha)\sin\delta\,\frac{z}{r_1},
\end{equation}
so the ``unit'' normal is not a unit vector. The error is maximised when
$\sin(2\alpha)\sin\delta=1$, giving
$\lVert\nhat\rVert-1=\sqrt{1+z/r_1}-1$, i.e. \SI{0.995}{\percent} at
$z/r_1=0.02$ and \SI{4.40}{\percent} at $z/r_1=0.09$. It vanishes identically at
$\alpha=0$, so no face-on result is affected.

\subsection{Differentiating the sail acceleration}

Write $k=\beta(1-\mu)/r_1^2$. From \eqref{eq:sail},
\begin{equation}
  \bm a = k\Bigl[\cos^3\!\alpha\,\rhat
        + \cos^2\!\alpha\sin\alpha\cos\delta\,\phat
        + \cos^2\!\alpha\sin\alpha\sin\delta\,\qhat\Bigr] .
\end{equation}
Differentiating with respect to the clock angle,
\begin{equation}
  \frac{\partial\bm a}{\partial\delta}
  = k\cos^2\!\alpha\sin\alpha\bigl[-\sin\delta\,\phat+\cos\delta\,\qhat\bigr],
  \label{eq:dadelta}
\end{equation}
which carries an overall factor $\sin\alpha$ and therefore
\begin{equation}
  \boxed{\;\left.\frac{\partial\bm a}{\partial\delta}\right|_{\alpha=0}=\bm0\;}
\end{equation}
identically, for every $\delta_0$. Geometrically: at $\alpha=0$ the normal is
parallel to $\rhat$, and the clock angle only rotates the normal about
itself, which does nothing.

For the cone angle, using
$\mathrm{d}(\cos^2\!\alpha\sin\alpha)/\mathrm{d}\alpha
=\cos^3\!\alpha-2\cos\alpha\sin^2\!\alpha$ and
$\mathrm{d}(\cos^3\!\alpha)/\mathrm{d}\alpha=-3\cos^2\!\alpha\sin\alpha$,
\begin{equation}
  \frac{\partial\bm a}{\partial\alpha}
  = k\Bigl[-3\cos^2\!\alpha\sin\alpha\,\rhat
    + \bigl(\cos^3\!\alpha-2\cos\alpha\sin^2\!\alpha\bigr)
      \bigl(\cos\delta\,\phat+\sin\delta\,\qhat\bigr)\Bigr] .
  \label{eq:dadalpha}
\end{equation}
At $\alpha=0$ this reduces to $k(\cos\delta_0\,\phat+\sin\delta_0\,\qhat)$: a
single, purely transverse direction fixed by the constant $\delta_0$. The radial
authority vanishes because
$\mathrm{d}(\cos^3\!\alpha)/\mathrm{d}\alpha\to0$ as $\alpha\to0$.

\subsection{Why rank of \texorpdfstring{$B$}{B} is the wrong test}

A face-on sail supplies one input direction, and it is transverse. One might
conclude the radial direction is unreachable. That is wrong, because the
Coriolis term in \eqref{eq:inplane} couples $\xi$ and $\eta$: a transverse push
generates radial motion within one dynamical time. The correct test is the
Kalman rank condition on $[B,\,MB,\,M^2B,\dots,M^5B]$ for the $6\times6$ system
matrix $M$. Applying it at the $\beta=0.05$ equilibrium ($c_2=1.409194$) gives
rank $4/6$ at $\alpha_0=0,\delta_0=0$ --- the in-plane four-state block is fully
controllable (rank $4/4$), and it is the \emph{out-of-plane} pair that is lost,
because $z$ decouples at linear order and only $\partial\bm a/\partial\delta$
could have reached it.

\subsection{The optimal cone angle}

The transverse thrust component is $\propto\cos^2\!\alpha\sin\alpha$. Maximising,
\begin{equation}
  \frac{\mathrm{d}}{\mathrm{d}\alpha}\bigl(\cos^2\!\alpha\sin\alpha\bigr)
  = \cos\alpha\bigl(\cos^2\!\alpha-2\sin^2\!\alpha\bigr) = 0
  \;\Longrightarrow\;
  \tan^2\!\alpha=\tfrac12 ,
\end{equation}
so
\begin{equation}
  \alpha_0^\star=\arctan\bigl(1/\sqrt2\bigr)=35.264^\circ,
  \qquad
  \cos^2\!\alpha_0^\star=\frac{1}{1+\tan^2\!\alpha}=\frac23 .
\end{equation}
This is McInnes' classical optimal sail cone angle; recovering it from the
control Jacobian is a benchmark, not a new result.

What is new is the conditioning there. At $\alpha_0^\star$,
$\cos^2\!\alpha-2\sin^2\!\alpha=\tfrac23-\tfrac23=0$, so the transverse part of
\eqref{eq:dadalpha} vanishes and
\begin{equation}
  \frac{\partial\bm a}{\partial\alpha}
  = -3k\cos^2\!\alpha\sin\alpha\,\rhat
  = -3k\cdot\tfrac23\cdot\tfrac{1}{\sqrt3}\,\rhat
  = -\frac{2k}{\sqrt3}\,\rhat ,
\end{equation}
purely radial, while from \eqref{eq:dadelta}
\begin{equation}
  \left\lVert\frac{\partial\bm a}{\partial\delta}\right\rVert
  = k\cos^2\!\alpha\sin\alpha = k\cdot\tfrac23\cdot\tfrac1{\sqrt3}
  = \frac{2k}{3\sqrt3} .
\end{equation}
The two columns are orthogonal, so they \emph{are} the singular vectors, and
\begin{equation}
  \frac{\sigma_2}{\sigma_1}
  = \frac{2k/(3\sqrt3)}{2k/\sqrt3} = \frac13
\end{equation}
exactly, independent of $\beta$ and $\mu$.

% =====================================================================
\section{Every residual, explained}\label{sec:resid}
% =====================================================================

This section exists because a reader meeting ``\num{2.1e-16}'' in an abstract
cannot tell whether it is a physical result, a convergence tolerance, or noise.
Each of the manuscript's small numbers is a different kind of quantity.

\subsection{Floating-point background}

All computation is IEEE-754 double precision. The machine epsilon is
$\varepsilon=2^{-52}=\num{2.220446e-16}$: the gap between $1$ and the next
representable number. For a value $x$, the gap to its neighbour --- one
\emph{unit in the last place}, or ulp --- is approximately $x\varepsilon$. A
disagreement of a few ulp between two computations of the same quantity means
they agree as well as double precision permits, and says nothing about the
mathematics.

\subsection{Table 2: the closed form against root-finding}

The column $\lvert\Delta\rvert$ is $\lvert\beta_{\mathrm{closed}} -
\beta_{\mathrm{brentq}}\rvert$, where the first evaluates
\eqref{eq:closedform} directly and the second solves $s(\beta)-1=0$ by Brent's
method through the numerically computed equilibrium. The two share no code path
beyond the balance itself, so agreement exercises the whole chain.

\begin{table}[t]
  \centering
  \caption{The Table-2 residuals expressed in units in the last place. All four
  are at the level of the root-finder's convergence, not of a disagreement in
  the mathematics.}
  \label{tab:ulp}
  \begin{tabular}{lcccc}
    \toprule
    system & $\beta_{\mathrm{crit}}$ & $\lvert\Delta\rvert$ &
    $\mathrm{ulp}(\beta)$ & $\lvert\Delta\rvert$ in ulp \\
    \midrule
    Sun--Mercury & \num{0.010961525} & \num{3.1e-16} & \num{1.7e-18} & \num{179} \\
    Sun--Earth   & \num{0.028646456} & \num{2.1e-16} & \num{3.5e-18} & \num{60} \\
    Sun--Jupiter & \num{0.187186696} & \num{2.8e-16} & \num{2.8e-17} & \num{10} \\
    Earth--Moon  & \num{0.406934946} & \num{5.0e-16} & \num{5.6e-17} & \num{9} \\
    \bottomrule
  \end{tabular}
\end{table}

Table~\ref{tab:ulp} shows the pattern: the residual is roughly constant in
absolute terms ($\sim\num{3e-16}$) rather than in relative terms, so it grows in
ulp as $\beta$ shrinks. That is the signature of an \emph{absolute} convergence
tolerance in the root-finder, which is exactly what \texttt{brentq} is
configured with. Nothing here is physics.

\subsection{The 61-point sweep}

The manuscript states that over $\mu\in[10^{-7},10^{-2}]$ sampled at 61
logarithmically spaced points, the two computations never differ by more than
\num{1.9e-15}. The worst case occurs at $\mu=\num{2.1544e-4}$, where
$\beta=0.116303$ and the residual is \num{1.8874e-15}, or $136$ ulp. The
Sun--Earth value \num{2.1e-16} quoted separately is that one system's residual
and is \emph{not} the sweep maximum --- an earlier draft conflated the two.

\subsection{The frequency scan}\label{sec:scanresid}

The manuscript reports that a $4\times10^5$-point scan places the minimum of
$\nu/\omega$ at $c_2=1.60000000$ with value $0.942809041582$, ``agreeing with
\eqref{eq:cstar} and \eqref{eq:bound} to \num{8.1e-10} and \num{1.1e-16}
respectively''. These two numbers mean very different things and the first is
misleading.

The scan is \texttt{linspace(1+1e-9, 4.2, 400001)}, a uniform grid of step
\begin{equation}
  h = \frac{4.2-(1+10^{-9})}{400000} = \num{8.0e-6} .
\end{equation}
A grid search can locate a minimum no better than about $h/2=\num{4e-6}$ in
general. The reported \num{8.125e-10} is the distance from $8/5$ to the
\emph{nearest grid point}, which happens to be small because of how the offset
$10^{-9}$ and the step interact --- an arithmetic accident of this particular
\texttt{linspace}, not a measure of resolving power. Had the grid been offset
differently, that figure would have been up to $\num{4e-6}$.

The second number is meaningful. At a stationary point the function is
quadratically flat, so a location error $\epsilon$ produces a value error
$\sim\tfrac12R''\epsilon^2$. With $\epsilon\sim10^{-9}$ the value error is far
below $\varepsilon$, which is why the \emph{value} agrees to \num{1.1e-16}
--- one ulp --- while the \emph{location} does not. The honest summary is:
the scan confirms the value of the bound to machine precision and is consistent
with the location $c_2=8/5$ to within its grid resolution of $\num{4e-6}$.

\subsection{The Jacobi spreads}\label{sec:jacobiresid}

``Spread'' means $\max-\min$ of the quantity sampled along one full period of a
closed orbit. On the first atlas member ($\beta=0.001$, $A_z=\num{5.0415e-4}$):
\begin{itemize}
\item $C_{\mathrm{sail}}$ spread $\approx\num{1.8e-15}$ --- this is
      \emph{integrator noise}. $C_{\mathrm{sail}}$ is an exact invariant, so its
      variation measures only the accumulated truncation and rounding error of
      the ODE solver. It is tolerance-dependent: re-integrating at
      \texttt{rtol=atol=1e-13} gives \num{1.33e-15} instead. Any value at the
      $10^{-15}$ level means ``conserved to machine precision''.
\item $C_{\mathrm{grav}}$ spread $=\num{5.787e-6}$ --- this is \emph{physical}.
      By \eqref{eq:cgrav} the two differ by $2\beta(1-\mu)/r_1$, which varies as
      $r_1$ varies along the orbit. It is nine orders of magnitude larger and is
      not a numerical artefact. It is the reason the sail potential, not the
      gravitational one, must be used.
\end{itemize}
Note that over an exactly closed orbit $r_1$ returns to its initial value, so
$C_{\mathrm{grav}}$ returns to its initial value too: the failure mode is not a
secular drift but a dependence on \emph{where} on the orbit the constant is
sampled. Sampling always at the same crossing makes the wrong quantity look
smooth, which is what concealed the error rather than exposing it.

\subsection{Other small numbers}

\begin{description}
\item[\num{3.07e-5}, the $L_3$ miss.] At $\mu=\tfrac12$ the classical $L_3$
reaches $c_2=1.569787$, giving $\nu/\omega$ larger than the bound
\eqref{eq:bound} by \num{3.07e-5}. This is a genuine dynamical near-miss, not a
numerical residual: $L_3$ occupies precisely the interval in which the extremum
$c_2^\star=8/5$ lies, and falls just short of reaching it.

\item[$<10^{-13}$, the $A\equiv c_2$ identity.] Two algebraically identical
expressions evaluated by different code paths. The residual is accumulated
rounding over a handful of operations on numbers of order unity.

\item[\num{4.5e-4} and \num{5.2e-2}, the singular values.] These are physical:
the second and first singular values of the transverse block of
$B_{\mathrm{sail}}$ at $\alpha_0=0.5^\circ$, in nd acceleration units. Their
ratio, $115{:}1$, is the conditioning penalty for operating near the singular
nominal. It is not a tolerance.

\item[$\Delta_z$, the branch discriminant.]\label{sec:deltaz} Defined as
$(\lvert z_0\rvert-\lvert z(T/2)\rvert)/(\lvert z_0\rvert+\lvert z(T/2)\rvert)$,
this is a dimensionless asymmetry, identically zero on the vertical-Lyapunov
branch and $O(0.1)$ on a genuine halo. From \texttt{halo\_atlas.csv} the
smallest value in each family falls from $0.109$ at $\beta=0.001$ to $0.0898$
at $\beta=0.040$ and $0.0125$ at $\beta=0.050$. It remains a valid
discriminator across the band --- the competing branch sits at zero --- but at
the upper end it separates the two by one order of magnitude rather than two.
\end{description}

% =====================================================================
\section{Physical scales for Sun--Earth}\label{sec:scales}
% =====================================================================

\begin{table}[h]
  \centering
  \caption{Conversions used throughout, $\mu=\num{3.003e-6}$.}
  \begin{tabular}{ll}
    \toprule
    quantity & value \\
    \midrule
    1 nd length              & \SI{1.495978707e8}{\kilo\metre} \\
    1 nd time                & \SI{58.1324}{\day} \\
    one full revolution      & $2\pi$ nd $=$ \SI{365.25636}{\day} \\
    Hill radius $r_H$        & \num{0.0100033} nd $=$ \SI{1496477}{\kilo\metre} \\
    classical $L_1$ standoff & \num{0.0099699} nd $=$ \SI{1491472}{\kilo\metre} \\
    parity standoff          & \num{0.0144273} nd $=$ \SI{2158294}{\kilo\metre} \\
    critical sail loading $\sigma^\ast$ & \SI{1.5311}{\gram\per\metre\squared} \\
    \bottomrule
  \end{tabular}
\end{table}

% =====================================================================
\section{Reproducing every number}\label{sec:repro}
% =====================================================================

The repository is public at
\url{https://github.com/Bishwaswarup/SOLAR_SAIL}. From a clean clone:
\begin{verbatim}
  pip install -r requirements.txt
  python main.py verify     # regression suite + 175 manuscript checks
  python main.py results    # regenerates results.txt
  python main.py atlas      # re-walks the 12 halo families (slow)
  python main.py figures    # regenerates the paper figures
\end{verbatim}
\texttt{paper/verify\_numbers.py} is the authority on manuscript agreement. It
recomputes each claim from the \texttt{src/} modules and compares it against the
literal typed in \texttt{main.tex}; Tables 2, 3 and 4 are parsed cell by cell
out of the manuscript source, so a table cannot silently disagree with the
formula it tabulates. A drifted number fails there rather than in review.

\end{document}
